基于持续同调的时间序列变点检验可利用滑动窗口提取局部结构特征，但窗口共享观测引起的依赖可能使随机置换参考分布失真。本文考察固定持续图向量表示与最大扫描所构成检验的有限样本校准。首先，以高斯窗口均值构造反例，说明原始观测独立同分布仍不足以保证派生窗口序列的交换性。随后，在固定窗口数与固定原始观测长度两种设计下，对 28 个零假设配置比较无约束窗口级随机置换与依赖乘子自助法，并以均值--协方差特征作为非拓扑对照。在名义显著性水平 0.05 下，对每种表示，20 个重叠配置均表现为实质反保守，持续同调表示的经验拒绝率最高为 0.9756。依赖乘子法在全部匹配配置中降低拒绝率，但对每种表示，8 个非重叠配置均表现为实质保守；原序列依赖对步长对比的修饰方向随生成过程与表示变化。低阶矩匹配的圆形到 8 字形曲线备择下，两类表示的尺寸校正功效分别为 1.0000 和 0.0620；整体平移负对照则呈相反结果，与距离型持续同调的平移不变性一致。结果表明，当前实现的校准表现具有明确的条件依赖性，降低拒绝率不能等同于恢复名义校准，表示的信号敏感性与检验的误报控制应分别评价。

\noindent\textbf{关键词}：持续同调；变点检测；重叠滑动窗口；交换性；有限样本校准。
